\subsection{Weights}\label{sga3:XXI:6.6}

\begin{definition}\label{sga3:XXI:6.6.1}
For a root datum \(\cR\), put
\[
 \Lambda(\cR)=
 \{x\in\VR(R)\mid(\alpha^\ast,x)\in\mathbb Z
   \text{ for every }\alpha^\ast\in R^\ast\}.
\]
The elements of \(\Lambda(\cR)\) are called the \emph{weights} of
\(\cR\). The weights of \(\cR^\ast\) are called the \emph{coweights}
of \(\cR\).
\end{definition}

One has
\[
 \Gamma_0(R)\subset\Lambda(\cR),
\]
and \(\Lambda(\cR)\) is stable under \(W(\cR)\).

\begin{lemma}\label{sga3:XXI:6.6.2}
The bilinear map
\[
 V^\ast\times V\longrightarrow\mathbb Q
\]
induces a duality
\[
 \Gamma_0(R^\ast)\times\Lambda(\cR)\longrightarrow\mathbb Z.
\]
\end{lemma}

\begin{proof}
Immediate.
\end{proof}

\begin{corollary}\label{sga3:XXI:6.6.3}
Let
\[
 \Delta^\ast=(\alpha_1^\ast,\ldots,\alpha_n^\ast)
\]
be a system of simple coroots, and let \(p_i\in\VR(R)\),
\(i=1,\ldots,n\), be defined by
\[
 (\alpha_i^\ast,p_j)=\delta_{ij}.
\]
Thus
\[
 s_{\alpha_i}(p_i)=p_i-\alpha_i,
 \qquad
 s_{\alpha_i}(p_j)=p_j\quad(i\ne j).
\footnote{Caution: when the datum is not reduced, the \(\alpha_i\)
need not form a system of simple roots.}
\]
Then \(\Lambda(\cR)\) is the free abelian group generated by the
\(p_i\).
\end{corollary}

The \(p_i\) are called the \emph{fundamental weights} corresponding to
the system of simple coroots \(\Delta^\ast\).

\begin{corollary}\label{sga3:XXI:6.6.4}
For every \(\alpha^\ast\in\Delta^\ast\),
\[
 \left(\alpha^\ast,\sum_i p_i\right)=1.
\]
Consequently,
\[
 \sum_i p_i=\rho_{R_+},
 \qquad
 R_+=\mathscr P\bigl(\operatorname{ind}(\Delta)\bigr),
\]
in the notation of \ref{sga3:XXI:3.5.1}.
\end{corollary}

\begin{corollary}\label{sga3:XXI:6.6.5}
For every \(x\in\VR(R)\),
\[
 x=\sum_i(\alpha_i^\ast,x)p_i.
\]
\end{corollary}

Since
\[
 R^\ast\subset\Gamma_0(R^\ast),
 \qquad
 R\subset\Lambda(\cR),
\]
the quadruple
\[
 \bigl(\Lambda(\cR),\Gamma_0(R^\ast),R,R^\ast\bigr)
\]
is a root datum.

\begin{corollary}\label{sga3:XXI:6.6.6}
The canonical morphism
\[
 \Gamma_0(R^\ast)\longrightarrow M^\ast
\]
is the transpose of
\[
 x\longmapsto\sum_i(\alpha_i^\ast,x)p_i.
\]
This map defines a morphism of root data, and there is a commutative
diagram
\begin{center}
\begin{tikzcd}[column sep=huge,row sep=large]
\cR \arrow[rr] \arrow[dr]
  && \operatorname{sc}(\cR) \arrow[dl,"\wr"'] \\
& \bigl(\Lambda(\cR),\Gamma_0(R^\ast),R,R^\ast\bigr) &
\end{tikzcd}
\end{center}
\end{corollary}
